\chapter{Introduction}
\label{sec:introduction}

Understanding what neural networks \textit{compute}, not only what they \textit{output},
is the central problem of mechanistic interpretability. A model that correctly solves arithmetic problems or chains
reasoning steps coherently need not be implementing anything resembling the algorithm a human would use.
Whether it does, where in the network that computation lives, and whether it can be manipulated predictably
are empirical questions that evaluation benchmarks cannot answer.

\section{Mechanistic Interpretability}

As neural networks scale from toy models to systems executed in consequential and complex settings,
the gap between performance and understanding becomes untenable. Thus, a model achieving near-human accuracy on reasoning benchmarks
cannot be straightforwardly understood. A system that has learned a spurious shortcut is indistinguishable from one that has learned the intended
algorithm until it fails at the margin.
The interpretability methods that accompanied earlier machine learning, such as saliency maps,
attribution scores, and probing classifiers, operate in the input or output space and offer
correlational accounts but are not mechanistic, as they describe what the model attends to, not what it computes.

Evidence for the tractability of mechanistic questions accumulated independently. Work on convolutional vision networks
identified low-level features, including edge detectors, curve analysers, and frequency filters,
that arise independently of architecture and training seed~\citep{olah2020zoom}. Analysis of transformer attention heads
revealed heads implementing specific, describable operations, such as induction heads completing
repeated token sequences and positional heads copying information from fixed offsets~\citep{elhage2021mathematical}.
A direct demonstration of such mechanisms is the grokking phenomenon~\citep{power2022grokking}, which shows that small models
trained on modular arithmetic initially memorise their training data and later undergo a phase transition,
developing compressed algorithms that generalise across all held-out operand pairs for the trained modulus.
The transition is invisible to accuracy metrics evaluated on the training set. However, Fourier decomposition
of the learned weight matrices reveals that the converged model implements modular arithmetic through
specific trigonometric identities encoded in the residual stream, an algorithm with no surface
similarity to the training objective~\citep{nanda2023progress}.

Further work has confirmed that small transformers trained on arithmetic tasks converge on specific,
inspectable algorithms, independent of surface statistical regularities in the
data~\cite{zhong2023clock, quirke2024understanding, quirkeneo2025understanding, musat2024clustering}.

The same interpretability considerations apply to large language models (LLMs), yet at a scale that makes manual
circuit analysis infeasible. As LLMs are trained on vastly more diverse data and without explicit algorithmic
curricula, they nonetheless achieve high accuracy on structured reasoning tasks. However, whether this reflects stable, localisable internal computation or brittle surface associations
cannot be determined from evaluation alone. This distinction matters wherever model outputs yield real consequences.

Anthropic's \textit{On the Biology of a Large Language Model}~\cite{lindsey2025biology} proposes
a reproducible methodology for approaching interpretability on a large language model (LLM).
Combining sparse autoencoders (SAEs) with cross-layer transcoders (CLTs) and backward-pass attribution,
it constructs dependency graphs that trace information flow from input tokens through intermediate
features to output logits. The framework forms the methodological foundation of this thesis, which
examines how structured reasoning is implemented within large instruction-tuned models and what
the geometry of those representations reveals about the nature of computation.

\section{Interpretability Methods}
\label{sec:interp-methods}

Let $\mathcal{M}$ be a transformer with $L$ layers, residual stream dimension $d$, and MLP
sublayers $\mathrm{MLP}^{(0)}, \ldots, \mathrm{MLP}^{(L-1)}$.  All methods employed in this
thesis operate on the residual stream of $\mathcal{M}$.  Let $\mathbf{h}^{(\ell)}_t \in \mathbb{R}^d$
denote the residual-stream vector at layer $\ell \in \{0, \ldots, L{-}1\}$ and token position $t$.
Each layer contributes an additive update,
\begin{equation}
    \mathbf{h}^{(\ell+1)}_t = \mathbf{h}^{(\ell)}_t
      + \mathrm{Attn}^{(\ell)}\!\left(\mathbf{H}^{(\ell)}\right)_t
      + \mathrm{MLP}^{(\ell)}\!\left(\mathbf{h}^{(\ell)}_t\right),
    \label{eq:residual}
\end{equation}
where $\mathbf{H}^{(\ell)}$ denotes the full sequence of residual-stream vectors at layer $\ell$.
This additive structure makes the residual stream a natural object for linear analysis, as each
layer writes into a shared representation that accumulates across depth.

Mechanistic analysis replaces each MLP sublayer $\mathrm{MLP}^{(\ell)}$ with a pretrained
\textit{transcoder} $\mathcal{T}_\ell$, a sparse encoder--decoder whose weights are fit offline
to reconstruct the MLP output while decomposing it into $K$ interpretable
directions~\citep{bricken2023monosemanticity, cunningham2023sparse}.  For the input
$\mathbf{x} \in \mathbb{R}^d$ to the MLP sublayer at layer $\ell$, the transcoder computes
pre-activations
\begin{equation}
  z_{\ell,k}(\mathbf{x})
    = \bigl(\mathbf{W}_\mathrm{enc}^{(\ell)}\,\mathbf{x}
           + \mathbf{b}_\mathrm{enc}^{(\ell)}\bigr)_k,
  \label{eq:preact}
\end{equation}
so that feature $k$ carries activation
$a_{\ell,k}(\mathbf{x}) = \mathrm{ReLU}(z_{\ell,k}(\mathbf{x}))$,
and the reconstructed MLP output is
\begin{equation}
  \hat{f}_\ell(\mathbf{x})
    = \mathbf{W}_\mathrm{dec}^{(\ell)}\,\mathbf{a}_\ell(\mathbf{x})
    + \mathbf{b}_\mathrm{dec}^{(\ell)},
  \label{eq:recon}
\end{equation}
where $\mathbf{W}_\mathrm{dec}^{(\ell)} \in \mathbb{R}^{d \times K}$ and
$\mathbf{a}_\ell(\mathbf{x})$ is the vector of all feature activations at layer $\ell$.
Features are indexed throughout by the double subscript $(\ell, k)$, with $\ell$ specifying
the layer and $k$ the index within that layer's dictionary.  Sparsity ensures that at most a
small active set contributes for any given input, making each feature's role inspectable in
isolation.

Attribution scores are computed on the \textit{linearised} model, in which attention outputs
are detached and treated as constants, so the residual stream propagates solely through the
transcoder reconstruction path.  This linearisation is the necessary condition for the scores
to be additive and interpretable as influence flows through the network.  Let $T$ denote the
position of the final input token, and let $y_j = \mathbf{v}_j^\top \mathbf{h}_T$ be the
logit for output token $j$, where $\mathbf{h}_T$ is the pre-unembedding hidden state and
$\mathbf{v}_j$ is the demeaned unembedding direction for token $j$.  The attribution score
of feature $k$ at position $t$ in layer $\ell$ is
\begin{equation}
    \alpha_{\ell,k,t}
      = a_{\ell,k,t}
        \cdot \frac{\partial y_j}{\partial a_{\ell,k,t}},
    \label{eq:attribution}
\end{equation}
the product of the feature activation with the gradient of the target logit with respect to
that activation~\citep{lindsey2025biology}.  Collecting all such scores over a given prompt
yields a directed weighted graph, the \textit{attribution graph}, in which nodes are
(layer, position, feature) triples and edges carry the influence $\alpha_{\ell,k,t}$ from
earlier features to later ones.  The graph is pruned by retaining only those nodes and edges
that account for a fixed fraction of the total downstream influence on $y_j$, controlled by
a node threshold $\tau_\mathrm{node}$ and an edge threshold $\tau_\mathrm{edge}$.

Activation patching establishes causal claims~\citep{conmy2023automated}. Given a clean prompt
$x$ and a counterfactual prompt $x'$, patching replaces the residual stream at layer $\ell$
and position $t$,
\begin{equation}
    \tilde{\mathbf{h}}^{(\ell)}_t = \mathbf{h}'^{(\ell)}_t,
    \label{eq:patching}
\end{equation}
and measures the resulting shift in output logit,
$\Delta_j = y_j(\tilde{\mathbf{h}}) - y_j(\mathbf{h})$.  A non-zero $\Delta_j$ confirms
that the patched representation is causally sufficient to alter the prediction.  Constrained
patching restricts the replacement to a targeted subset of positions or components, enabling
finer localisation of causal dependence.

Steering vectors are extracted by contrasting two prompt sets, $D_+$ where a concept is
present and $D_-$ where it is absent, and computing their mean difference at a designated
anchor position $t^*$,
\begin{equation}
    \delta^{(\ell)} = \mathbb{E}_{x \in D_+}\!\left[\mathbf{h}^{(\ell)}_{t^*(x)}\right]
                 - \mathbb{E}_{x \in D_-}\!\left[\mathbf{h}^{(\ell)}_{t^*(x)}\right].
    \label{eq:delta}
\end{equation}
Under the linear representation hypothesis~\citep{park2023linear}, $\delta^{(\ell)}$
approximates the axis along which the concept is encoded at depth $\ell$.  Injecting
$\alpha\,\delta^{(\ell)}$ into the residual stream at inference time provides a continuous
lever on the concept representation without modifying model weights.  The norm profile
$\lVert\delta^{(\ell)}\rVert$ across layers is the primary signal used for concept
localisation throughout this thesis.  In combination, these four methods support mechanistic
claims about how a model computes, not only whether it succeeds on a given input.

\section{This Thesis}

We take the open-access model \texttt{Qwen3-4B} as our subject. The primary contribution is a
systematic concept localisation study across a diverse set of concepts spanning three structural categories:
modular arithmetic (carry propagation in integer addition, GCD divisibility, and residue class membership),
relational reasoning (transitive ordering and negation scope), and physical and causal plausibility
(energy conservation and causal direction). The central hypothesis is that modular concepts,
whose outcomes are determined by discrete thresholds or periodic relations, produce qualitatively different localisation signatures from relational and
physical concepts, whose outcomes depend on monotone comparison or world knowledge.

The rest of this dissertation is structured as follows. Chapter~\ref{sec:related}
covers the relevant background. ......
